\section{Conclusion}\label{sec:conclusion}

We presented SoundCode, the first asynchronous design for
verifier-guided code generation. The system is a producer-consumer
decoding loop in which a compiler verifier runs in a
background task at each depth-zero structural boundary, the producer
LM streams tokens without blocking, and a consumer rolls back the
buffer to the last verifier-confirmed checkpoint only on a returned
blocking diagnostic. The same rollback algorithm underlies the
synchronous ROCODE-style baseline \citep{jiang2024rocode}; the
sync-versus-async axis is isolated at fixed verifier, fixed model,
and fixed benchmark.

\paragraph{Where async pays.}
The win grows with verifier latency. Empirically, on Rust ---
where \code{cargo check} is $\sim$$6$--$18\times$ a per-token
decoding step in our setup --- async overlap recovers most of the verifier
wall-time ($1.59\times$ speedup; async wins on $116$ of $156$
paired problems). On C++, where each
$g{+}{+}\;\code{-fsyntax-only}$ invocation amortizes to
$\sim$$0.6$\,s of per-boundary cost, the win is
larger still ($2.95\times$; async wins on $159$ of $161$ paired
problems) --- as the cost-spectrum analysis in
Section~\ref{sec:when-async} predicts. pass@1 is statistically
indistinguishable between the schedulings on both languages
(\S\ref{sec:main-results}); the result is a latency result. We expect the win to grow
further on \code{javac}-supervised Java
($\sim$$180\,\mathrm{ms}$) and to be largest on
whole-crate borrow / trait analysis (up to seconds), neither of
which we evaluate here but both of which the same architecture
supports without modification.

\paragraph{Limitations.}
The matched-budget caveat from \citet{olausson2024selfrepair}
applies: any speedup is only meaningful if the verifier itself is a
strictly-stronger oracle than the generator on the disputed
fragment. We argue that a compiler \emph{is} such an oracle on the
decidable fragment of syntax, type, and lifetime correctness ---
unlike a learned classifier \citep{wang2025semguard}, the compiler
has no false-positive rate on errors that fall inside its formal
spec. But this argument is restricted to that decidable fragment:
the compiler tells us nothing about functional correctness, and any
test-pass gain SoundCode delivers over vanilla decoding has to come
from compilable-but-wrong code being a small fraction of all
generated code. We do not address this gap; SemGuard
\citep{wang2025semguard} is the natural complement that targets
exactly the semantic-but-compilable failure mode.

The current evaluation is also single-model and single-seed, and
we accordingly make no accuracy claims: none of the pass@1
differences between arms is statistically significant
(\S\ref{sec:main-results}), and the contribution of this paper is
confined to wall-clock latency. Five of the nine arms in the
design space --- including the AsyncLSP arm, the only one that
would test LSP-based (rather than compiler-based) verification ---
were not run (\S\ref{sec:setup}), so every claim here is about
compiler supervision. Multi-seed confidence intervals, the
AsyncLSP arm, and a faithful SemGuard retraining on Rust
diagnostics are the obvious next experiments.

\paragraph{Future work.}
Six extensions are immediate. (0) \emph{The AsyncLSP arm.} The
rust-analyzer driver is implemented but unevaluated; measuring it
is the missing experiment that would extend this paper's claim
from compiler supervision to LSP supervision.
(1) \emph{Oscillation-detection rollback.} On HumanEval-140 (\code{fix\_spaces}) under our
\textsf{SyncCargo} arm, the model hit the 60\,s problem budget after
emitting $66$ syntactically-valid repetition statements that all
re-introduced the same compile error in the same place; each one was
individually cleared by the boundary detector and re-rejected by the
verifier. A policy that detects ``$N$ consecutive rollbacks landing
within an offset window of each other'' and escalates to a coarser
target (halve the buffer; restart from the prompt with a penalty
mask) would have caught this pathology in tens of milliseconds. The
case is concrete evidence that the next paper-sized contribution on
this architecture is a \emph{penalty / oscillation} mechanism
layered on top of the rollback policy. (2) \emph{Faithful SemGuard
comparison.} SemGuard is the strongest learned-verifier baseline;
adapting it from Python/Java to Rust requires retraining the 1.3B
classifier on a curated Rust diff corpus, work currently in
progress. (3) \emph{Larger benchmarks.} MBPP-rs and LiveCodeBench
\citep{jain2024livecodebench} would extend the evaluation from
single-function HumanEval to multi-function and contamination-free
problems. (4) \emph{Beam-search and DFS rollback policies.} Our
unifying $\pi$ abstraction admits multi-branch policies that we
have not evaluated; the natural reference is
LATS \citep{zhou2024lats}. (5) \emph{Composition with cheap
front-end masks.} DOMINO \citep{beurerkellner2024domino} and
Type-Constrained \citep{mundler2025typeconstrained} are
near-zero-overhead syntactic prevention layers; composing them with
SoundCode's expensive semantic verifier-with-rollback back-end
should be Pareto-better than either alone.

\paragraph{Artifacts.}
A live web demo and the reference implementation are open-source.
Demo:
\href{https://soundcode-demo.psixyzt.com}%
     {\code{soundcode-demo.psixyzt.com}}.
Source:
\href{https://github.com/leobwang/soundcode}%
     {\code{github.com/leobwang/soundcode}}
\citep{soundcode2025github}.
